\chapter{Software para modelización, análisis, procesamiento y síntesis de señales con Python}

\section{Introducción a Python, NumPy y Matplotlib}

Python es un lenguaje de programación de propósito general, pero al equiparlo con librerías científicas como \texttt{NumPy} (para cálculo numérico y arreglos) y \texttt{Matplotlib} (para graficación), se convierte en un entorno inmensamente poderoso, similar a Matlab/Octave, pero de código completamente abierto.

En este apunte, mostramos algunas de las operaciones y comandos básicos para la operación interactiva utilizando Python.

\subsection{La línea de comandos y el entorno}

\begin{enumerate}
	
	\item Hoy en día, la forma más común de interactuar con Python para cálculo científico es a través de \textbf{Jupyter Notebooks} o consolas interactivas como \textbf{IPython}. En estos entornos, escribimos el código en "celdas" o líneas de comandos marcadas con \verb-In []:-.
	
	\item Para tener acceso a las funciones matemáticas y gráficos, siempre comenzaremos nuestros scripts importando las librerías fundamentales:
	\begin{verbatim}
		import numpy as np
		import matplotlib.pyplot as plt
	\end{verbatim}
	
	\item Para probar operaciones básicas podemos escribir:
	\begin{verbatim}
		2 + 2
		3 - 5
		3 / 1.4
		2 * np.pi / 4
		(1 + 2j) / (1 - 2j)
	\end{verbatim}
	\textbf{Nota:} En Python, la unidad imaginaria se representa con \texttt{j} en lugar de \texttt{i}, y siempre debe llevar un número delante (ej. \texttt{1j} en lugar de solo \texttt{j}).
	
	\item Si utilizamos un entorno tradicional, podemos guardar nuestras órdenes en un archivo de texto con extensión \texttt{.py} (un "script") y ejecutarlo completo desde la terminal con \texttt{python mi\_script.py}.
\end{enumerate}

\subsection{El sistema de ayuda}

\begin{enumerate}
	\item Para ver cómo se usa un comando en Python, utilizamos la función integrada \texttt{help()}:
	\begin{verbatim}
		help(np.cos)
	\end{verbatim}
	\item Si utilizamos IPython o Jupyter, podemos colocar un signo de interrogación al final de la función para abrir una ventana de ayuda rápida:
	\begin{verbatim}
		np.cos?
	\end{verbatim}
\end{enumerate}

\subsection{Valores escalares}

\begin{enumerate}
	\item Como hemos visto, operamos con números reales y complejos, también llamados escalares.
	\item Probemos los siguientes comandos:
	\begin{verbatim}
		(1 + 2j) * (1 - 2j)
		45 * (56 + 1j) - 34 * (16 - 32j)
	\end{verbatim}
\end{enumerate}

\subsection{Valores matriciales o arreglos}

\begin{enumerate}
	\item NumPy trabaja con un objeto fundamental llamado \texttt{ndarray} (n-dimensional array). Para crear vectores y matrices usamos la función \texttt{np.array()}.
	\item Podemos crear un vector fila escribiendo:
	\begin{verbatim}
		c = np.array([1 + 2j, 1 - 3j])
	\end{verbatim}
	
	\item Podemos crear un vector columna anidando listas (una lista de listas, donde cada sublista es una fila):
	\begin{verbatim}
		e = np.array([[1 + 2j], 
		[1 - 3j]])
	\end{verbatim}
\end{enumerate}

\subsection{Metáfora del entorno de trabajo}

Podemos suponer, a modo de metáfora, que disponemos de un pizarrón donde escribimos las cuentas que vamos haciendo, línea por línea.

\subsection{Variables}

\begin{enumerate}
	
	\item Python permite asignar nombres a los diferentes valores mediante el operador \texttt{=}. El símbolo \texttt{\#} se usa para indicar que aquello escrito a su derecha es un comentario y no será ejecutado por el intérprete.
	
	\item Pruebe ingresar los siguientes valores:
	\begin{verbatim}
		a = 5                  # escalar real
		b = 5 + 10j            # escalar complejo
		c = np.array([1, 2, 3]) # vector de 1D
		d = np.array([[1], [2], [3]]) # vector columna (matriz 3x1)
		f = np.array([[1, 2, 3], 
		[4, 5, 6], 
		[7, 8, 9]]) # una matriz 3 x 3
	\end{verbatim}
\end{enumerate}

\subsection{Operaciones con números complejos}

\begin{enumerate}
	\item A continuación vemos algunas de las operaciones que pueden realizarse con números complejos en NumPy. Pruebe ingresar:
	\begin{verbatim}
		x = 3 + 4j
		np.real(x)  # 3.0
		np.imag(x)  # 4.0
		np.conj(x)  # 3.0 - 4.0j
		np.abs(x)   # 5.0 (Módulo)
		np.angle(x) # 0.9273 (Fase en radianes)
	\end{verbatim}
\end{enumerate}

\subsection{Raíces y coeficientes de un polinomio}

\begin{enumerate}
	\item Supongamos que tenemos un polinomio general:
	\begin{equation}
		b_{n}x^{n}+b_{n-1}x^{n-1}+...+b_{0}=0.
	\end{equation}
	Para hallar las raíces del polinomio podemos colocar sus coeficientes en un arreglo ordenado (de mayor a menor grado) y usar la función \texttt{np.roots}:
	\begin{verbatim}
		p = [bn, ..., b0]
		np.roots(p)
	\end{verbatim}
	
	\item Supongamos que tenemos que hallar las raíces del polinomio $x^{2}+2x+3=0$. Escribamos entonces:
	\begin{verbatim}
		p = [1, 2, 3]
		np.roots(p)
	\end{verbatim}
	
	\item Supongamos ahora que conocemos las raíces y necesitamos hallar los coeficientes del polinomio originario. Podemos usar la función \texttt{np.poly}:
	\begin{verbatim}
		r = [-1.0, 2.7, -3.1, -3.1]
		np.poly(r)
	\end{verbatim}
	
	\item ¿Cuál le parece que es el resultado de la siguiente orden?
	\begin{verbatim}
		np.poly(np.roots([1, 2, 3, 4, 5]))
	\end{verbatim}
	
\end{enumerate}

\subsection{Generación de vectores y matrices}

\begin{enumerate}
	\item NumPy permite la generación de secuencias de valores. La función equivalente al operador dos puntos (\texttt{:}) es \texttt{np.arange(inicio, fin, paso)}.
	\item \textbf{¡CUIDADO! DIFERENCIA CRÍTICA:} En Python, el valor de \texttt{fin} \textbf{NUNCA} se incluye en el arreglo resultante.
	\begin{verbatim}
		t = np.arange(1, 2.1, 0.1)
	\end{verbatim}
	Se genera el vector \texttt{[1.0, 1.1, ..., 2.0]}. Tuvimos que usar \texttt{2.1} como fin para que incluyera el \texttt{2.0}.
	
	\item Una forma mucho más segura para números con decimales es \texttt{np.linspace(inicio, fin, cantidad\_de\_puntos)}, que \textbf{sí} incluye el final:
	\begin{verbatim}
		t = np.linspace(1, 2, 11) # 11 puntos entre 1 y 2
	\end{verbatim}
	
	\item Para generar matrices llenas de ceros o unos podemos usar las funciones \texttt{np.zeros} y \texttt{np.ones}, pasándoles una tupla con las dimensiones:
	\begin{verbatim}
		Z = np.zeros((3, 3))
		U = np.ones((2, 4))
	\end{verbatim}
\end{enumerate}

\subsection{Operaciones sobre vectores y matrices, primera parte}

\begin{enumerate}
	\item Para averiguar las dimensiones de un arreglo en NumPy usamos el atributo \texttt{.shape}:
	\begin{verbatim}
		t = np.linspace(1, 2, 11)
		print(t.shape) # Imprime (11,)
	\end{verbatim}
	
	\item Las operaciones algebraicas de suma, resta y multiplicación por escalares se realizan elemento a elemento de forma nativa:
	\begin{verbatim}
		t + 2
		t - 2
		t * 2
		np.cos(t)
		np.sin(2 * t - 1)
	\end{verbatim}
	
	\item \textbf{¡CUIDADO! OTRA DIFERENCIA CRÍTICA:} A diferencia de otros lenguajes matemáticos, en NumPy el operador de multiplicación \texttt{*} y exponenciación \texttt{**} actúan \textbf{elemento a elemento por defecto}. No necesita poner un punto (\texttt{.}).
	\begin{verbatim}
		t * t   # Multiplicación elemento a elemento
		t ** 2  # Eleva al cuadrado cada elemento
	\end{verbatim}
	
	\item Si desea realizar un \textbf{producto matricial} verdadero (álgebra lineal), debe usar el operador \texttt{@} o la función \texttt{np.dot()}:
	\begin{verbatim}
		A = np.array([[1, 2], [3, 4]])
		B = np.array([[5, 6], [7, 8]])
		
		C = A * B  # Producto elemento a elemento
		D = A @ B  # PRODUCTO MATRICIAL REAL
	\end{verbatim}
\end{enumerate}

\subsection{Operaciones sobre vectores y matrices, indexación y extracción}

\begin{enumerate}
	\item Ingresemos un vector:
	\begin{verbatim}
		v = np.array([10, 20, 30, 40, 50, 60, 70])
	\end{verbatim}
	
	\item \textbf{¡ÍNDICES DESDE CERO!:} En Python, el primer elemento está en la posición \texttt{0}, no en la \texttt{1}.
	\begin{verbatim}
		v[0] # Extrae el primer elemento (10)
		v[2] # Extrae el TERCER elemento (30)
	\end{verbatim}
	
	\item También podemos extraer una porción (slice) del vector escribiendo \verb-v[i:j]-. En Python, esto extrae desde el índice \texttt{i} \textbf{hasta el índice j-1} (el límite superior no se incluye):
	\begin{verbatim}
		v[1:4] # Extrae los elementos en índices 1, 2 y 3: [20, 30, 40]
	\end{verbatim}
	
	\item Podemos acceder al último elemento dinámicamente usando índices negativos:
	\begin{verbatim}
		v[-1] # Extrae el 70
		v[1:] # Extrae desde el índice 1 hasta el final
		v[:-1] # Extrae desde el principio hasta el penúltimo
	\end{verbatim}
\end{enumerate}

\subsection{Álgebra lineal}

\begin{enumerate}
	\item El sub-módulo \texttt{np.linalg} contiene la librería de álgebra lineal. 
	\begin{center}
		\begin{tabular}{ll}
			\textbf{Comando} & \textbf{Resultado} \\ \hline
			\texttt{np.linalg.det(A)} & Determinante \\
			\texttt{np.linalg.eig(A)} & Autovalores y autovectores \\
			\texttt{np.linalg.inv(A)} & Inversa matricial \\
			\texttt{np.linalg.svd(A)} & Descomposición SVD \\
			\texttt{np.linalg.cond(A)} & Número de condición \\
		\end{tabular}
	\end{center}
	
	\item Supongamos un sistema de 2 ecuaciones lineales $Ax=b$:
	\begin{verbatim}
		A = np.array([[1, 2], [2, 1]])
		b = np.array([[1], [2]])
	\end{verbatim}
	
	\item El método recomendado y optimizado numérica y computacionalmente para resolver el sistema es usar la función \texttt{solve()}:
	\begin{verbatim}
		x = np.linalg.solve(A, b)
	\end{verbatim}
	
	\item Verifique que el resultado es correcto calculando el vector de error (residuo):
	\begin{verbatim}
		e = (A @ x) - b
		np.transpose(e) @ e   # Suma de los errores al cuadrado (cercana a 0)
	\end{verbatim}
\end{enumerate}

\subsection{Graficación de funciones 2D con Matplotlib}

\begin{enumerate}
	\item Para graficar funciones usamos el módulo \texttt{pyplot} de Matplotlib. Escriba:
	\begin{verbatim}
		t = np.arange(0, 2*np.pi, 0.1)
		y = np.cos(2 * t - 1)
		
		plt.figure(1)
		plt.plot(t, y)
		
		plt.figure(2)
		plt.stem(t, y)
		plt.show() # Este comando abre las ventanas en pantalla
	\end{verbatim}
	
	\item \texttt{plot} dibuja una señal de tiempo continuo uniendo los puntos con líneas, mientras que \texttt{stem} grafica la secuencia de valores discretos explícitos.
	
	\item Estos comandos admiten un parámetro de formato idéntico al de Matlab (\texttt{'r*:'}):
	\begin{itemize}
		\item \textbf{Colores:} \texttt{y}, \texttt{m}, \texttt{c}, \texttt{r}, \texttt{g}, \texttt{b}, \texttt{w}, \texttt{k}.
		\item \textbf{Marcadores:} \texttt{.}, \texttt{o}, \texttt{x}, \texttt{*}, \texttt{s}.
		\item \textbf{Líneas:} \texttt{-}, \texttt{:}, \texttt{-.}, \texttt{--}.
	\end{itemize}
	
	\item Observe los comandos complementarios para enriquecer el gráfico:
	\begin{verbatim}
		plt.figure()
		plt.plot(t, np.cos(t), 'm:o')
		plt.grid(True)
		plt.xlabel('Tiempo (s)')
		plt.ylabel('Amplitud (V)')
		plt.title('Función periódica')
		plt.show()
	\end{verbatim}
	
	\item El comando \texttt{subplot} permite colocar múltiples gráficos en la misma ventana (filas, columnas, índice activo):
	\begin{verbatim}
		plt.figure(figsize=(10,8)) # Ajusta el tamaño de la ventana
		plt.subplot(2, 1, 1); plt.plot(t, np.cos(t))
		plt.subplot(2, 1, 2); plt.stem(t, np.cos(t))
		plt.show()
	\end{verbatim}
\end{enumerate}

\subsection{Cómo guardar gráficos en archivos}

\begin{enumerate}
	\item Una vez dibujada la figura (y \textbf{antes} de llamar a \texttt{plt.show()}), podemos guardar el gráfico en alta calidad (PNG o PDF) utilizando:
	\begin{verbatim}
		plt.savefig('mi_grafico.png', dpi=300)
	\end{verbatim}
\end{enumerate}

\subsection{Graficación de funciones en tres dimensiones}

\begin{enumerate}
	\item Es posible graficar curvas y superficies en 3D importando herramientas específicas e indicándole a la figura que la proyección es 3D:
	\begin{verbatim}
		fig = plt.figure()
		ax = fig.add_subplot(projection='3d')
		
		# Ejemplo rápido
		X, Y = np.meshgrid(np.linspace(-5, 5, 50), np.linspace(-5, 5, 50))
		Z = np.sin(np.sqrt(X**2 + Y**2))
		ax.plot_surface(X, Y, Z, cmap='viridis')
		plt.show()
	\end{verbatim}
\end{enumerate}



\subsection{Escritura y salida de texto formateado}

\begin{enumerate}
	\item Con la función nativa \texttt{print} podemos mostrar texto y valores. En Python moderno, lo mejor son los "f-strings" (cadenas formateadas), anteponiendo una \texttt{f} y colocando las variables entre llaves \texttt{\{\}}:
	\begin{verbatim}
		variable = 42
		print(f"El resultado del cálculo es {variable}")
		print(f"El valor de pi es aproximadamente {np.pi:.4f}")
	\end{verbatim}
	El modificador \texttt{:.4f} indica que se muestren 4 decimales.
\end{enumerate}

\subsection{Estructuras de Control: Iteración e identación}

\begin{enumerate}
	\item \textbf{¡CUIDADO! IDENTACIÓN:} En Python no se usa la palabra \texttt{end} para terminar un ciclo o condición. El bloque de código que pertenece al ciclo se define **únicamente** por su margen izquierdo (sangría o identación).
	
	\item Bucle \texttt{for} (la función \texttt{range(1, 9)} genera enteros del 1 al 8):
	\begin{verbatim}
		for x in range(1, 9):
		y = 2 * x - 1
		print(f"x = {x}; y = {y}")
		# Aquí, al quitar la sangría, el bucle terminó.
	\end{verbatim}
	
	\item Alternativamente, usando un bucle \texttt{while}:
	\begin{verbatim}
		x = 1
		while x <= 8:
		y = 2 * x - 1
		print(f"x = {x}; y = {y}")
		x = x + 1 
	\end{verbatim}
\end{enumerate}

\subsection{Toma de decisiones}

\begin{enumerate}
	\item El resto de la división entera se calcula con el operador \texttt{\%}.
	\item Podemos condicionar la ejecución empleando \texttt{if}, \texttt{elif}, y \texttt{else}:
	\begin{verbatim}
		for x in range(1, 16):
		if x % 2 == 0:
		print(f"{x} es par")
		elif x % 3 == 0:
		print(f"{x} es impar y divisible por 3")
		else:
		print(f"{x} es impar")
	\end{verbatim}
\end{enumerate}

\section{Ejemplo 1: Integración gráfica}

En esta sección repasaremos los conceptos escribiendo un pequeño script para:
\begin{enumerate}
	\item Generar un vector eje de abscisas de 51 números.
	\item Generar datos aleatorios asociados a ese eje.
	\item Graficar con etiquetas correspondientes.
\end{enumerate}

Para graficar los números generados, escriba el siguiente código y ejecútelo:
\begin{verbatim}
	x = np.arange(1, 52)         # Vector de 1 a 51
	y = np.random.rand(len(x))   # Ruido uniforme entre 0 y 1
	
	plt.plot(x, y, 'b-*')
	plt.title('Señal estocástica generada por computadora')
	plt.xlabel('Índice de muestra (n)')
	plt.ylabel('Amplitud aleatoria')
	plt.grid(True)
	plt.show()
\end{verbatim}

\section{Ejemplo 2: Interacción con archivos de texto (I/O)}

En Python, la lectura y escritura de matrices numéricas está extremadamente simplificada gracias a NumPy, evitando los ineficientes bucles anidados.

\begin{enumerate}
	\item \textbf{Escritura directa:} Supongamos que generamos una matriz de datos \texttt{A} de $3 \times 10$:
	\begin{verbatim}
		# Generamos datos de prueba
		filas = np.arange(1, 4).reshape(3, 1) # Vector columna
		cols = np.arange(1, 11).reshape(1, 10) # Vector fila
		A = filas + cols  # Broadcasting mágico de NumPy genera matriz 3x10
		
		# Guardamos en archivo en una sola línea
		np.savetxt('matriz_datos.txt', A, fmt='%d')
	\end{verbatim}
	
	\item \textbf{Lectura directa:}
	\begin{verbatim}
		# Leemos el archivo y lo convertimos a matriz de enteros automáticamente
		matriz_leida = np.loadtxt('matriz_datos.txt', dtype=int)
		print(matriz_leida.shape) # Verificará que es (3, 10)
	\end{verbatim}
\end{enumerate}

\section{Ejemplo 3: Vectorización eficiente y protección matemática}

En el siguiente algoritmo procesaremos datos evitando el uso ineficiente de ciclos \texttt{for} y protegiéndonos de cálculos indefinidos (divisiones por cero).

\begin{enumerate}
	\item La indexación lógica \verb-(np.abs(vector) < eps)- nos permite encontrar y enmascarar valores peligrosamente cercanos a cero.
	\item \texttt{eps} es la constante de "precisión de máquina", accesible vía \texttt{np.finfo(float).eps}.
	
	\begin{verbatim}
		# Supongamos que tenemos una matriz_leida de Nx3 dimensiones
		a = matriz_leida[:, 0]  # Extrae TODA la columna 0 (la primera)
		b = matriz_leida[:, 1]
		c = matriz_leida[:, 2]
		
		eps = np.finfo(float).eps
		
		# Protegemos el divisor 'c' de valer exactamente 0
		ud = c + (np.abs(c) < eps) * eps 
		vd = (a - b)
		vd = vd + (np.abs(vd) < eps) * eps
		
		# Cálculos vectorizados: División punto a punto automática en NumPy
		u = (a + b) / ud  
		v = c / vd        
		
		print(u)
		print(v)
	\end{verbatim}
\end{enumerate}